% ============================================================
%  Chapter 2 — Literature Review
% ============================================================
\chapter{Literature Review}
\label{chap:lit}

\section{Overview of Electricity Price Forecasting}

The electricity price forecasting (EPF) literature has expanded rapidly since
market liberalisation began in the 1990s. \citet{weron2014} provides the
definitive survey, classifying approaches into five families: (1) multi-agent
models, (2) fundamental models, (3) reduced-form stochastic models,
(4) statistical models, and (5) computational intelligence models. The survey
covers over two decades of research and remains the standard reference for
the field, identifying that no single method dominates across all markets and
horizons, and that the most practically useful improvements typically come from
better feature engineering rather than more complex model architectures.
More recently, \citet{nowotarski2018} document the shift towards probabilistic
forecasting and machine learning methods, arguing that ensemble combination
of diverse forecasting approaches consistently outperforms individual models.

The field has converged on several stylised facts about electricity price
dynamics: (1) prices exhibit multi-frequency seasonality (intraday, weekly,
annual); (2) the distribution is heavily right-skewed with non-trivial negative
mass; (3) extreme price spikes are driven by fundamental supply-demand
imbalances rather than by speculative dynamics; and (4) the price-setting
technology shifts across hours and seasons, making the relationship between
fundamentals and prices inherently non-linear. These characteristics motivate
the use of non-linear ensemble methods capable of capturing threshold effects
and interaction terms.

For the purposes of this thesis, we focus on short-term \textit{point}
forecasting at the day-ahead horizon, which remains the most practically
relevant setting for physical market participants who must submit bids to the
EPEX SPOT auction by 12:00 CET each day. The review is structured around three
questions: (1) what methods have been proposed and evaluated? (2) what role
does weather play in these models? and (3) how should models be formally
compared?

\section{Statistical and Econometric Approaches}

Early work on day-ahead price forecasting relied on time-series models adapted
from financial econometrics. Autoregressive models — in particular ARIMA and
its seasonal extensions (SARIMA) — capture the strong autocorrelation structure
of electricity prices but struggle to model the non-linear interactions between
price and its fundamental drivers. GARCH-family models address volatility
clustering \citep{weron2014} but are ill-suited to the extreme price spikes
that characterise electricity markets, where the distribution of returns has
much heavier tails than financial asset returns, including a non-trivial mass
of negative observations \citep{fanone2013}.

The mean-reversion and jump-diffusion framework, introduced in the commodity
pricing literature, is better able to capture the spike-and-reversion behaviour
of electricity prices. However, these structural models typically require
numerical methods for calibration and are difficult to extend to include
exogenous weather or fundamental variables, limiting their usefulness when
the research objective is to quantify the marginal contribution of individual
predictors.

Factor models and regularised regression represent an important middle ground.
\citet{maciejowska2016} apply Factor Quantile Regression Averaging to
day-ahead prices across European markets, combining multiple linear
specifications through optimal weights. Critically, \citet{uniejewski2019}
demonstrate that automated variable selection using LASSO regularisation applied
to a rich set of price lags and calendar variables achieves competitive
performance relative to more complex non-linear models, even on markets with
high renewable penetration. This result — that well-specified linear models
with carefully chosen price lags can be surprisingly competitive — is directly
relevant to our finding that weather adds negligible marginal value once price
lags and fundamental variables are included. The common thread is that the
predictive information in electricity prices is highly concentrated in recent
price history: any additional variable, whether weather-based or
fundamental-based, must provide information \textit{not already present} in
recent prices to justify its inclusion.

\section{Machine Learning Approaches}

The application of machine learning to electricity price forecasting has
accelerated since the mid-2010s, driven by the increasing availability of
high-resolution data and cheap computational resources. \citet{aggarwal2009}
provide an early review, while \citet{lago2018} offer an empirical comparison
of deep learning architectures (LSTM, CNN) against traditional benchmarks on
European day-ahead markets, finding that deep learning improves on ARIMA but
not consistently over well-tuned tree ensembles.

The most systematic evaluation of machine learning methods for EPF is provided
by \citet{lago2021}, who benchmark 27 algorithms across 12 European electricity
markets using a standardised open-source evaluation framework. Their key
finding is that no algorithm consistently dominates: the best-performing method
varies by market and time period, and the margin of improvement from complex
deep learning over simpler tree-based methods is often statistically
insignificant once proper Diebold-Mariano testing is applied. For the French
market specifically, \citet{lago2021} report benchmark MAE values in the range
of 10--25 EUR/MWh under stable market conditions, which provides the reference
frame against which our results (MAE $\approx$ 16.94 EUR/MWh in May 2024 --
April 2025) should be interpreted. The authors explicitly recommend that EPF
papers adopt DM testing as standard practice for model comparison, a
recommendation we follow throughout Chapter~\ref{chap:model}.

\textbf{Tree-based ensemble methods} — Random Forests \citep{breiman2001} and
Gradient Boosted Trees (including XGBoost, \citealt{chen2016}) — have proven
particularly well-suited to the EPF problem for several reasons. First, they
handle non-linear interactions between weather, calendar, and fundamental
variables without requiring manual feature transformations. Second, they are
robust to outliers and near-zero prices, which are frequent in liberalised
markets and create difficulties for methods that rely on symmetric loss
functions. Third, tree ensembles provide built-in feature importance measures
(Mean Decrease in Impurity for Random Forests) that offer interpretable
insights into price drivers, a property highly valued by market participants
who must justify their models to risk committees. \citet{tschora2022} confirm in a
recent comparative study on European day-ahead markets that gradient boosted trees
remain competitive with neural network architectures on EPF tasks when sample sizes
are moderate ($n < 100{,}000$ hourly observations), as in our setting.

\citet{ziel2018} compare univariate and multivariate forecasting frameworks on
European day-ahead markets and find that multivariate models incorporating
load, generation, and cross-border flows outperform univariate price-only
models, supporting the use of fundamental variables in our feature matrix.
A consistent finding across the ML literature is that price lags dominate
feature importance rankings when included in the feature matrix \citep{lago2021,
uniejewski2019}. This is consistent with the strong autocorrelation structure
of electricity prices and the difficulty of improving over recent price history
as a predictor. Our results — where the 24-hour price lag accounts for
approximately 22\% of total MDI importance — are in line with this consensus.

\section{Weather as a Driver of Electricity Prices}

The causal pathway from weather to electricity prices is well-established in
the fundamental literature. Temperature affects demand through heating (in
winter) and cooling (in summer). Wind speed determines the output of wind
generators, which can set the marginal cost to zero when in excess, driving
prices downward. Solar irradiance compresses midday prices during
high-penetration periods, a phenomenon known as the ``solar merit order
effect''. Precipitation affects hydro reservoir levels, with droughts leading
to higher thermal generation and higher prices.

\citet{staffell2016} document the strong correlation between wind speed and
wind generation output (approximately cubic relationship), which underlies our
\textit{wind power proxy} feature.

In the French market specifically, the thermosensitivity is unusually high by
European standards. \citet{rte2023} report that each degree Celsius below
17\textdegree C increases peak demand by approximately 2.4 GW — a consequence
of the high penetration of electric heating (over 30\% of French households).
This motivates the use of 17\textdegree C as the Heating Degree Day threshold
throughout this thesis, consistent with RTE's own operational methodology and
higher than the 15\textdegree C threshold used in Germany.

However, the literature has increasingly recognised that weather effects on
electricity prices may be \textit{indirect} rather than direct — transmitted
through other observable variables. If electricity prices are primarily set by
the marginal cost of gas-fired generation, and if gas prices themselves respond
to pan-European heating demand, then a model including TTF gas prices may
already capture the meteorological signal without explicit weather features.
This indirect transmission channel is the foundation of our information
redundancy hypothesis.

The 2021--2023 European energy crisis provides a natural experiment on this
hypothesis. The simultaneous withdrawal of Russian pipeline gas supply and
France's historically low nuclear availability (reaching $\approx$40\% in
mid-2022 due to corrosion-related maintenance) created a regime in which gas
prices became the dominant driver of electricity prices across Europe. In such
a regime, the direct weather channel — cold temperatures increasing electricity
demand directly — may be more visible above the gas price noise than in normal
market conditions. Testing whether weather features are more informative during
the 2022 crisis than in the 2024--25 normalisation period is a central
contribution of this thesis (Section~\ref{sec:validation2022}).

\section{Diebold-Mariano Testing in Energy Forecasting}

Comparing forecast accuracy across models requires careful statistical
treatment. Informal comparisons based on point estimates of MAE or RMSE can
lead to misleading conclusions when differences are not statistically
significant. The Diebold-Mariano (DM) test \citep{diebold1995} provides a
formal framework for testing the null hypothesis of equal predictive accuracy
between two competing forecasts. The test statistic is based on the difference
in loss functions (here, absolute errors) and accounts for autocorrelation in
the loss differential series via a Newey-West long-run variance estimator
\citep{newbold1993}.

For small samples, \citet{harvey1997} propose a correction factor that adjusts
the DM statistic to better approximate a $t$-distribution, reducing the risk
of over-rejection of the null. We apply this Harvey-Leybourne-Newbold (HLN)
correction throughout our analysis, following best practice in the EPF
literature \citep{weron2014}. \citet{lago2021} explicitly recommend the
DM test with HLN correction as the standard for EPF model comparisons, noting
that many published results that appear significant on point metrics alone are
not statistically significant once proper testing is applied. Adopting this
standard is particularly important for our weather ablation study: the
difference in MAE between our models B and C is only 0.01 EUR/MWh — a
difference that appears negligible on inspection and is confirmed as
statistically insignificant by the DM test ($p = 0.572$).

\section{Gap in the Literature}

Despite the growing body of work on machine learning for electricity price
forecasting, several gaps remain relevant to this thesis:

\begin{enumerate}
  \item \textbf{French market specificity:} Most empirical studies focus on the
  German, Spanish, or Nordic markets, or treat European markets as homogeneous.
  The French market's combination of dominant nuclear share ($\approx$70\% of
  generation in normal conditions), unusually high thermosensitivity, and
  interconnection with multiple neighbours creates a distinct price formation
  mechanism. The 2022 crisis, during which France's nuclear availability fell
  to an unprecedented 40\%, represents a structurally unique episode whose
  impact on the weather-price relationship has not been formally studied in the
  EPF literature.

  \item \textbf{Formal ablation of weather features:} While weather variables
  are routinely included in EPF models, their marginal contribution is rarely
  isolated using a statistically formal test. Most studies either exclude
  weather entirely or include it without testing whether its addition is
  statistically significant relative to a fundamental-only baseline. We address
  this gap with a Diebold-Mariano ablation study that directly tests whether
  ERA5 meteorological variables improve on a model that already includes load
  forecasts, fuel prices, and nuclear availability.

  \item \textbf{Regime-dependence of the weather contribution:} No study, to
  our knowledge, explicitly tests whether the marginal contribution of weather
  variables varies across market regimes (crisis vs.\ normal). Our design,
  which evaluates Models B vs C on both the 2024--25 stable period and the 2022
  crisis period using the same DM testing framework, directly addresses this gap
  and produces our most novel empirical finding.

  \item \textbf{Economic value evaluation:} Most studies evaluate models on
  statistical metrics alone. We extend the analysis to an executable day-ahead
  trading backtest with realistic transaction costs and a dead-band filter
  (Chapter~\ref{chap:backtest}), linking statistical accuracy to economic value
  and providing a practitioner-relevant assessment of model quality.
\end{enumerate}
